\section{Raft Internals}
\label{sec:RaftInternals}

This section is for developers changing the Raft machinery itself.

\subsection{Core state}

A Raft node maintains:

\begin{itemize}
\item current term,
\item voted-for candidate,
\item local role: follower, candidate or leader,
\item replicated log,
\item commit index and last-applied index,
\item known leader,
\item cluster membership configuration,
\item leader-side replication progress for followers,
\item read-lease and heartbeat timing state.
\end{itemize}

The safety rules are the usual Raft rules: terms are monotonic, a node votes at
most once per term, candidates must have sufficiently up-to-date logs, leaders
append entries in order, and committed entries must not be lost across leader
changes.

\subsection{Elections}

Election behavior is driven by timeout and vote responses:

\begin{enumerate}
\item A follower whose election deadline expires becomes a candidate.
\item The candidate increments term, votes for itself and sends vote requests.
\item A candidate becomes leader only after satisfying the configured majority rule.
\item Any node that observes a higher term steps down.
\end{enumerate}

Membership matters. In stable configuration, election success requires a
majority of current voters. In joint consensus, election success requires a
majority of current voters and a majority of next voters.

\subsection{AppendEntries and commit}

Leaders replicate log entries using AppendEntries. Followers reject requests
whose previous index/term does not match their log. Leaders backtrack until the
matching prefix is found, then append missing entries.

Commit advancement must respect both term and membership:

\begin{itemize}
\item The leader should only commit entries known replicated on a required majority.
\item Joint consensus requires both current and next majorities.
\item Entries from older terms require the normal Raft care: leadership should be
established through an entry in the current term before relying on older entries
for commitment.
\end{itemize}

\subsection{Read lease and read barrier}

Linearizable reads require proof that the leader is still authoritative. The
library uses a read-lease/read-barrier model:

\begin{itemize}
\item A leader records successful quorum contact through heartbeat/append responses.
\item While that evidence is fresh, it may serve local queries.
\item If evidence is stale, it must refresh a quorum heartbeat barrier.
\item If the barrier fails, the query returns retry/redirect behavior rather than
serving stale state.
\end{itemize}

The same membership rules apply here: joint consensus requires current and next
majorities for the barrier to be safe.

\subsection{Joint consensus}

Membership changes are replicated through internal log entries. A safe transition
uses two phases:

\begin{enumerate}
\item Enter joint consensus with current and next configurations active.
\item Finalize into the next stable configuration after the joint entry is committed.
\end{enumerate}

During joint consensus, majority calculations must be split, not flattened. A
single large majority across the union is not sufficient. The implementation must
evaluate current voters and next voters independently.

\subsection{Internal commands}

Internal Raft commands are log entries whose meaning is Raft-owned rather than
domain-owned. Membership changes are the main example. Internal commands must be
decoded and applied before or alongside domain state updates so that effective
configuration by log index remains correct.

Application developers should not create internal commands directly. Machinery
developers should keep internal command encoding versioned and compatible across
Java and C++.
